\section{Full coherent coefficient-tail pairs with a common first skeleton}
\label{sec:136-coherent-common-skeleton}

\paragraph{Archive and proved inputs.}
The uniform tests are cote136, physical pp.~117 and119,
footers116 and118, scans117 and119; the common-source question
is physical p.~121, footer120, scan121. We use the actual
common first skeleton \(F\) from AQ-136-FIRST-SKELETON and
the newly checked ARCH-136-TAIL-EXT-VANISHING. The latter
identifies the dual of every positive coefficient-tail resolution,
in every nonnegative weight shift, with one coefficient group
in degree zero, including after reduction by any integer.
No unproved Ext vanishing or inverse-limit exchange is used.

Write \(S=\Gamma_{\mathbb Z}(T)\). Let
\(f_j:A\to C_j\) be a finite family of strict maps, where
the \(C_j\) are admissible, the first skeletons of \(A\)
and all \(C_j\) are literally \(F\), and every \(f_j\)
is the identity there. Fix \(q\geq4\), a weight \(v>0\),
an integer \(N\geq1\), and elements in one simplex
\[
 z\in A_n^{v,q},\qquad b\in A_n^{v,q-1},\qquad
                         dz=0,\quad db=Nz.              \tag{136.CG1}
\]
The original weight \(v\) need not itself be in a target
comparison range. Take a sufficiently large coefficient tail
\(I_M\) and any graded free resolution
\[
 \cdots\longrightarrow P_1\longrightarrow P_0
                 \xrightarrow{\epsilon}I_M\longrightarrow0
                                                               \tag{136.CG2}
\]
whose generators all have weights at least \(M\).
Shift its weights by \(v\).

\paragraph{ARCH-136-SKELETON-COHERENT-MOORE.}
One can choose \(M\), using only the finitely many targets'
vertical and horizontal cutoffs through form degree \(q+1\)
and the endpoint-filling cutoff of \(F\), so that the
following hold uniformly in \(n\) and \(N\). For each
target there are \(S\)-linear maps
\[
 \begin{aligned}
 U_i^{(j)} &:P_i\langle v\rangle\longrightarrow
                         (C_j)_n^{*,q-1-i}
                              &&(0\leq i\leq q-1),\\
 V_i^{(j)} &:P_i\langle v\rangle\longrightarrow
                         (C_j)_n^{*,q-2-i}
                              &&(0\leq i\leq q-2),
 \end{aligned}                                             \tag{136.CG3}
\]
which give a full coherent Moore pair on all coefficient-tail
multiples of (136.CG1). Moreover all of these maps have
the same restrictions to
\(G=\operatorname{sk}_1\Delta[n]\), as cochains of the
common graph algebra \(F(G)\).

In particular the complete coherent presentation can be realized
by an actual simplicial strict PD extension under \(A\), over
all the \(C_j\), whose first skeleton is still literally
\(F\). It is rationally a relative free unit-disk extension
in each simplex. If \(A\) has the rational coefficient
comparison required by AQ-136-RELATIVE-H1, its unwanted
degree-zero and degree-one cohomology can subsequently be
repaired while preserving that first skeleton and all target
maps. There is no conclusion here about the other new integral
cohomology of the PD realization.

\begin{proof}
For a cochain complex \(C\), put
\[
 \mathcal H(C)^r=\bigoplus_{i=0}^r
        \operatorname{Hom}^0_{S,\mathrm{gr}}
                (P_i\langle v\rangle,C^{*,r-i}),
 \quad
 (\delta f)_i=df_i-(-1)^r f_{i-1}\partial_i.              \tag{136.CG4}
\]
Put \(K=\operatorname{Hom}^0_{S,\mathrm{gr}}
(P_\bullet\langle v\rangle,S)\) with this same total
convention. ARCH-136-TAIL-EXT-VANISHING gives the canonical
quasi-isomorphism \(S^v[0]\to K\). Since its terms are
products of torsion-free coefficient groups, its Moore cone
\[
 P_N(K)^r=K^r\oplus K^{r+1},\qquad
 d(x,y)=(\delta x-Ny,-\delta y)
\]
has just \(S^v/NS^v\) in degree zero.

Choose \(M+v\) above all the stated target bounds and the
endpoint bound for \(F\). The same first-quadrant comparison
used in ARCH-136-DERIVED-TAIL-COMPARISON, now evaluated in one
simplex, identifies the cohomology of
\(\mathcal H((C_j)_n)\) with \(K\) through total degree
\(q+1\). Indeed each column is a product of target complexes
in weights at least \(M+v\), and products are exact. Only
finitely many columns occur in any fixed total degree. The
canonical coefficient map supplies the comparison. It follows
from the cone exact sequence that
\[
 H^r(P_N(\mathcal H((C_j)_n)))=0
                         \quad(1\leq r\leq q).          \tag{136.CG5}
\]
No localization of an infinite Hom product occurs in this
argument.

We give the graph comparison explicitly, since it is the
additional issue. The common skeleton has \(F_0=S\), with
\(F_1\) a tensor product of free even zero-form disks.
It has a weight-preserving \(S\)-linear integral contraction
onto its coefficients. Endpoint surjectivity in all required
weights gives the exact graph-incidence sequence
\[
 0\longrightarrow F(G)\longrightarrow
       S^{V(G)}\oplus F_1^{E(G)}
       \longrightarrow (S\oplus S)^{E(G)}\longrightarrow0
                                                               \tag{136.CG6}
\]
after application of each Hom functor from the shifted free
\(P_i\). To be precise, the unqualified sequence may fail
surjectivity in small weights, but every free generator now
has weight at least the endpoint bound. Hom from that free
module is a product of exactly those surjective sequences;
its sequence is exact. The graph is finite.

The coefficient inclusion \(S\to F_1\) is a cochain
homotopy equivalence over \(S\), compatible with both
endpoint evaluations. Thus the homotopy fiber of the last
map in (136.CG6), after total Hom, is quasi-isomorphic to
the ordinary graph cochain complex with coefficients in
\(K\). One can see this without choosing simultaneous
endpoint retractions: compare the strict incidence diagram
formed with \(S\) in place of each \(F_1\) to (136.CG6).
The middle coefficient map is a quasi-isomorphism, the last
map is the identity, and the ordinary kernel of the actual
\(F_1\) incidence diagram computes its homotopy fiber because its map is
degreewise onto. The coefficient diagram's homotopy fiber
is the usual two-term graph complex, after eliminating one
copy of the edge coefficient. Taking the Moore cone and
using the concentration of \(P_N(K)\) proves
\[
 H^r(P_N(\mathcal H(F(G))))
              \cong H^r(G;S^v/NS^v),
       \qquad H^{r\geq2}(P_N(\mathcal H(F(G))))=0.
                                                               \tag{136.CG7}
\]
These identifications are induced by the coefficient maps.
All products and totalizations used above have exact products
inside terms and finite diagonals in any fixed total degree.

Let \(f_z\) and \(f_b\) have only their zeroth resolution
components, defined by
\(f_z(u)=\epsilon(u)f_j(z)\) and
\(f_b(u)=\epsilon(u)f_j(b)\). Then
\((f_b,f_z)\) is a closed element of
\(P_N(\mathcal H((C_j)_n))\) in degree \(q-1\).
By (136.CG5) it has a primitive \(y_j=(-V^{(j)},-U^{(j)})\)
in degree \(q-2\). Its equations are
\[
                       \delta U^{(j)}=f_z,\qquad
                       \delta V^{(j)}=NU^{(j)}-f_b.      \tag{136.CG8}
\]
All restrictions \(y_j|_G\) have the same differential.
By (136.CG7), their differences, which are cycles in degree
\(q-2\geq2\), are boundaries. Choose a cone primitive
for each difference on \(G\). Each component of such a
primitive extends to the corresponding target simplex:
the target row is acyclic, the relative normalized chains of
\(G\subset\Delta[n]\) are degreewise free, and the
coefficient parameters \(P_i\) are free, so the ordinary
row-extension argument applies independently to every basis
element. Subtracting the differential of that extension from
\(y_j\) leaves (136.CG8) unchanged and makes all graph
restrictions identical. This proves the asserted existence
with common graph values.

One can in fact arrange a particularly simple common graph
value. Edgewise contract the restrictions of \(z,b\) using
the \(S\)-linear contraction of \(F_1\). The forms have
degrees at least three. There are common graph cochains
\(c_G,t_G\) of degrees \(q-1,q-2\), with
\(dc_G=z|_G\), \(dt_G=Nc_G-b|_G\); they vanish on vertices.
Take the common zeroth components to be
\(\epsilon c_G,\epsilon t_G\), and all higher resolution
components zero. This satisfies (136.CG8) on the graph,
because \(\epsilon\partial_1=0\). The same exact-difference
argument arranges these specified restrictions in every target.

For completeness, write \(\eta=(-1)^q\). The full component
equations, including the bottom ones, are
\[
 \begin{aligned}
 dU_0&=f_z,&dU_i&=-\eta U_{i-1}\partial_i\quad(i\geq1),\\
 dV_0&=NU_0-f_b,&dV_i&=NU_i+\eta V_{i-1}\partial_i
                                                   \quad(i\geq1),\\
 U_{q-1}\partial_q&=0,&
 NU_{q-1}+\eta V_{q-2}\partial_{q-1}&=0.
 \end{aligned}                                             \tag{136.CG9}
\]
Introduce these generators on the representable simplex,
identify their graph values with the old elements just chosen
in \(F(G)=A(G)\), and impose the two bottom relations.
The free strict PD universal property gives all target maps.
The square of the differential and the differential stability
of the bottom relations follow from \(\partial^2=0\),
\(df_b=Nf_z\), and the two displayed signs. The full
identifying PD ideal is stable by the special differential
rule. The graph relations are stable because their values
obey the same equations. This is an actual strict extension,
with its first skeleton unchanged.

Pointwise its bottom parameter module is the signed version
of \(P+N^{-1}Q\) from ARCH-136-COHERENT-PD-REALIZATION;
integer saturation is proved in the same way. After
rationalization, eliminate the bottom \(U_{q-1}\) and use
\[
 Z_0=NU_0-f_b,\qquad
 Z_i=NU_i+\eta V_{i-1}\partial_i\quad(i\geq1).
\]
They are free coordinates with \(dV_i=Z_i\), \(dZ_i=0\).
The Euler contraction of the positive-arity part proves the
same rational relative equivalence as before, including for
arbitrary parameter ranks. AQ-136-RELATIVE-H1 is therefore
applicable under the hypothesis stated in the theorem.

Finally that repair keeps the first skeleton literally \(F\),
not merely its target images. In simplicial degrees zero and
one the source already is \(F\), whose normalized vertical
zero-to-one differential is an isomorphism onto one-cycles.
The replacement map on zero-form parameter modules used by
RH1 is consequently an isomorphism in those simplicial degrees.
Its PD pushout changes nothing there. The same is true at
every iteration. Coefficients survive by face retractions,
and all installed differential identities persist. This proves
the claimed strict-realization and low-repair assertions,
without an integral nonflat Poincare lemma.
\end{proof}

\paragraph{Actual degree-four coefficient relations.}
For the native first case, let \(a\) have form degree one,
put \(da=x\), \(z_0=\gamma_2(x)\), and \(e_0=ax\),
so \(de_0=2z_0\). For selected homogeneous \(s\in I_M\),
choose a lift \(e_s\in P_0\) and abbreviate
\(c_s=U_0(e_s)\), \(t_s=V_0(e_s)\). These depend on the
chosen lifts; no coefficient-linear splitting on \(I_M\)
is asserted. They obey
\[
 dc_s=s z_0,\qquad dt_s=2c_s-s e_0,
                 \qquad s\in I_M.
\]
Choose likewise a lift \(e_r\in P_0\) of \(r\in I_M\).
For a genuine coefficient relation \(fs=2r\), choose
\(h\in P_1\) with \(\partial h=f e_s-2e_r\).
The degree-three cycle
\(\zeta=f c_s-r e_0\), with \(2\zeta=d(f t_s)\), then
has the coherent primitive and its relation
\[
 \zeta=d(t_r-U_1(h)),\qquad
 2(t_r-U_1(h))-f t_s=d(-V_1(h)).                         \tag{136.CG10}
\]
These follow directly from (136.CG9) with \(q=4\).
With the simple common graph choice in the proof,
\(U_1,V_1\) vanish on \(G\), while
\(t_s|_G=s t_G\) and \(t_r|_G=r t_G\).
Thus these actual coefficient-syzygy pairs have compatible
graph restrictions; their degree-three graph residue is zero.
This resolves that particular compatibility condition, rather
than merely the abstract pointwise annihilator calculation.

\paragraph{Exact remaining nonlinear gate.}
The common graph lift works for every torsion pair of degree
\(q\geq4\), including non-native input pairs. For \(q=3\)
the same graph comparison leaves the ordinary group
\(H^1(G;S^v/NS^v)\); it does not vanish by the coefficient
Ext theorem. Moreover the finite strict realization reaches
form degree zero, has nonflat parameter modules, and can
create new integral cohomology beyond its linear parameter
complex. Formula (136.CG10) identifies and repairs the actual
linear relations produced by the initial native tail step;
it does not assert that all new degree-three classes of the
PD realization are of that form. Controlling those nonlinear
classes, while preserving all uniform admissibility bounds,
remains necessary for an integral admissible common source.
